% ============================================================
% CHAPTER 2 — REQUIREMENTS ANALYSIS
% ============================================================
\chapter{Requirements Analysis}

\section*{Introduction}
This chapter identifies the system actors, enumerates functional and non-functional requirements, presents the global use case diagram, and defines the system architecture.

% ---- 2.1 ACTORS ----
\section{Actors Identification}

The IGS Marketplace platform serves six distinct actor types organized in a hierarchical RBAC system:

\begin{table}[H]
\centering
\caption{System Actors and Permission Levels}
\begin{tabular}{L{2.5cm} C{2cm} C{2cm} L{7cm}}
\toprule
\textbf{Role}    & \textbf{Type}    & \textbf{Level} & \textbf{Key Permissions} \\
\midrule
super\_admin     & Internal         & 100            & Full system access, manage roles, system settings \\
admin            & Internal         & 80             & Manage users, assets, orders, view analytics \\
employee         & Internal         & 60             & Moderate content, view all analytics \\
intern           & Internal         & 40             & Browse catalog, request free assets, view tasks \\
seller           & External         & 20             & Upload assets, manage own listings, view analytics \\
member           & External         & 10             & Purchase assets, add to cart, write reviews \\
\bottomrule
\end{tabular}
\end{table}

% ---- 2.2 FUNCTIONAL REQUIREMENTS ----
\section{Functional Requirements}

\begin{table}[H]
\centering
\caption{Functional Requirements Summary}
\begin{tabularx}{\textwidth}{L{3.5cm} X}
\toprule
\textbf{Domain} & \textbf{Requirements} \\
\midrule
Authentication       & Register via email/password or Google OAuth2; JWT-based API authentication; profile management; admin role management. \\
\midrule
Asset Catalog        & Browse, filter (category, price, rating, free/paid), and search assets; 13 categories; interactive 3D/HDRI/texture viewer. \\
\midrule
Seller Upload        & Multi-format upload (GLB, FBX, OBJ, PNG, HDR\ldots); pending moderation workflow; edit request after approval; seller analytics dashboard. \\
\midrule
Shopping \& Orders   & Persistent cart; Stripe payment; order tracking; automatic library unlock; free asset claiming. \\
\midrule
Reviews \& Ratings   & Ownership-gated review submission; 1--5 star rating; AI sentiment analysis; automatic moderation flagging. \\
\midrule
Admin Panel          & Approve/reject assets with written reason; manage users and roles; platform-wide analytics; audit trail. \\
\midrule
Notifications        & Real-time and persistent in-app notifications; unread badge; mark as read. \\
\midrule
Intern Workflow      & Asset request with justification; voucher generation; supervised upload. \\
\midrule
Community            & P2P asset trading; community groups; seller role request flow. \\
\midrule
IGS Workspace        & Internal app: team management, task tracking, project ERP, real-time chat (Socket.io), role-based navigation. \\
\midrule
AI Services          & Personalized recommendation engine; marketplace chatbot (Mistral-7B); automated review sentiment analysis. \\
\bottomrule
\end{tabularx}
\end{table}

% ---- 2.3 NON-FUNCTIONAL REQUIREMENTS ----
\section{Non-Functional Requirements}

\begin{table}[H]
\centering
\caption{Non-Functional Requirements}
\begin{tabularx}{\textwidth}{L{3.5cm} X}
\toprule
\textbf{Category}        & \textbf{Requirements} \\
\midrule
\textbf{Performance}     & API response under 300 ms for standard queries; asset files streamed to cloud storage; queries indexed. \\
\textbf{Security}        & JWT expiration; Firebase token verification on all protected routes; RBAC middleware; parameterized queries; CORS; HTTPS in production. \\
\textbf{Scalability}     & Stateless REST API; cloud file storage (Supabase); connection pooling; horizontal scaling ready. \\
\textbf{Usability}       & Responsive UI for desktop and tablet; loading states and error messages on all user-facing operations. \\
\textbf{Reliability}     & Centralized error handling; database transactions for orders; audit logging for admin actions. \\
\textbf{Maintainability} & Routes → Controllers → Repositories separation; modular React components; consistent naming conventions. \\
\bottomrule
\end{tabularx}
\end{table}

% ---- 2.4 GLOBAL USE CASE ----
\section{Global Use Case Diagram}

\begin{figure}[H]
\centering
\resizebox{0.85\textwidth}{!}{% Global Use Case Diagram -- All Actors
\begin{tikzpicture}[
  actor/.style={rectangle, draw=isiBlue, fill=blue!10, rounded corners=3pt,
                minimum width=1.8cm, minimum height=0.65cm, align=center,
                font=\footnotesize\bfseries},
  uc/.style={ellipse, draw=isiBlue!60, fill=blue!4,
             minimum width=3.6cm, minimum height=0.65cm,
             align=center, font=\footnotesize, inner sep=1pt},
  extends/.style={draw=gray, dashed, ->, font=\tiny\itshape},
  assoc/.style={draw=gray!70, thin},
  boundary/.style={draw=isiBlue, very thick, fill=gray!3, rounded corners=6pt}]

  % ---- System Boundary ----
  \draw[boundary] (-4.2,-7.6) rectangle (4.2,5.2);
  \node[font=\footnotesize\bfseries\color{isiBlue}] at (0,5.0) {$\ll$system$\gg$ IGS Marketplace};

  % ---- Use Cases ----
  \node[uc] (browse)      at (0, 4.3)  {Browse Catalog};
  \node[uc] (search)      at (0, 3.4)  {Search \& Filter Assets};
  \node[uc] (viewdetail)  at (0, 2.5)  {View Asset Details (3D)};
  \node[uc] (cart)        at (0, 1.6)  {Manage Cart};
  \node[uc] (purchase)    at (0, 0.7)  {Purchase Asset (Stripe)};
  \node[uc] (download)    at (0,-0.2)  {Download / Library};
  \node[uc] (review)      at (0,-1.1)  {Write Review \& Rate};
  \node[uc] (upload)      at (0,-2.0)  {Upload Asset};
  \node[uc] (analytics)   at (0,-2.9)  {View Analytics};
  \node[uc] (moderate)    at (0,-3.8)  {Moderate Assets};
  \node[uc] (manageusers) at (0,-4.7)  {Manage Users};
  \node[uc] (reqasset)    at (0,-5.6)  {Request Asset (Intern)};
  \node[uc] (voucher)     at (0,-6.5)  {Redeem Voucher};
  \node[uc] (workspace)   at (0,-7.3)  {Access Workspace};

  % ---- Left Actors ----
  \node[actor] (visitor)  at (-7.2, 4.3)  {Visitor};
  \node[actor] (member)   at (-7.2, 1.0)  {Member};
  \node[actor] (seller)   at (-7.2,-2.4)  {Seller};
  \node[actor] (intern)   at (-7.2,-6.0)  {Intern};

  % ---- Right Actors ----
  \node[actor] (admin)    at (7.2,-3.3)   {Admin};
  \node[actor] (superadm) at (7.2,-6.5)   {Super Admin};
  \node[actor] (employee) at (7.2,-0.2)   {Employee};

  % ---- Associations ----
  % Visitor
  \draw[assoc] (visitor) -- (browse);
  \draw[assoc] (visitor) -- (search);
  \draw[assoc] (visitor) -- (viewdetail);

  % Member (inherits visitor)
  \draw[assoc] (member) -- (browse);
  \draw[assoc] (member) -- (cart);
  \draw[assoc] (member) -- (purchase);
  \draw[assoc] (member) -- (download);
  \draw[assoc] (member) -- (review);

  % Seller
  \draw[assoc] (seller) -- (upload);
  \draw[assoc] (seller) -- (analytics);
  \draw[assoc] (seller) -- (download);

  % Intern
  \draw[assoc] (intern) -- (browse);
  \draw[assoc] (intern) -- (reqasset);
  \draw[assoc] (intern) -- (voucher);
  \draw[assoc] (intern) -- (workspace);

  % Admin
  \draw[assoc] (admin) -- (moderate);
  \draw[assoc] (admin) -- (manageusers);
  \draw[assoc] (admin) -- (analytics);

  % Employee
  \draw[assoc] (employee) -- (moderate);
  \draw[assoc] (employee) -- (workspace);

  % Super Admin (extends Admin)
  \draw[assoc] (superadm) -- (manageusers);
  \draw[assoc] (superadm) -- (moderate);
  \draw[extends] (superadm) -- node[right, font=\tiny] {$\ll$extends$\gg$} (admin);

\end{tikzpicture}
}
\caption{Global Use Case Diagram}
\label{fig:use_case_global}
\end{figure}

% ---- 2.5 SYSTEM ARCHITECTURE ----
\section{System Architecture}

The IGS platform follows a \textbf{three-tier distributed architecture} separating presentation, business logic, and data access:

\begin{itemize}
  \item \textbf{Tier 1 — Presentation:} Two independent React + Vite SPAs (Marketplace on port 5173 and IGS Workspace on port 5174), communicating with the backend exclusively via REST and sharing a JWT stored in \texttt{localStorage}.
  \item \textbf{Tier 2 — Business Logic:} A single Express.js REST API (port 3001) handling authentication, authorization, domain logic, and external service orchestration (Stripe, Supabase, Firebase, AI).
  \item \textbf{Tier 3 — Data:} A PostgreSQL database (30+ tables) for relational persistence, and Supabase Storage for binary asset files.
\end{itemize}

Within Tier 2, the backend follows the \textbf{Repository pattern} — each entity has its own repository file encapsulating all SQL queries. Controllers never write raw SQL; they always go through a repository method. This decouples business logic from database access and makes the data layer independently testable.

\begin{figure}[H]
\centering
\resizebox{\textwidth}{!}{% Backend Layered Architecture (3-Tier + Repository Pattern)
\begin{tikzpicture}[
  layer/.style={rectangle, draw=isiBlue, fill=blue!6, rounded corners=4pt,
                minimum width=10cm, minimum height=1.1cm, align=center,
                font=\small\bfseries, text=isiBlue},
  module/.style={rectangle, draw=igsGreen!80, fill=green!5, rounded corners=3pt,
                 minimum width=2.6cm, minimum height=0.85cm, align=center,
                 font=\footnotesize, text=black},
  cross/.style={rectangle, draw=orange!80, fill=orange!8, rounded corners=3pt,
                minimum width=2.2cm, minimum height=0.75cm, align=center,
                font=\footnotesize, text=black},
  arr/.style={->, thick, draw=isiBlue},
  lbl/.style={font=\tiny\itshape, midway, left}]

  % ── Tier 1 : Presentation / Route Layer ──────────────────────────
  \node[layer] (routes) at (0, 6) {Tier 1 — Presentation / API Routes Layer};
  \node[module] at (-3.5, 5.1) {\texttt{routes/auth.js}};
  \node[module] at (-1.1, 5.1) {\texttt{routes/assets.js}};
  \node[module] at (1.3,  5.1) {\texttt{routes/orders.js}};
  \node[module] at (3.7,  5.1) {\texttt{routes/users.js}};

  % ── Middleware band ───────────────────────────────────────────────
  \node[draw=darkGray, fill=gray!8, rounded corners=3pt,
        minimum width=10cm, minimum height=0.65cm, align=center,
        font=\footnotesize\itshape, text=darkGray] at (0, 4.3)
        {Middleware: \texttt{authenticate} · \texttt{requirePermission} · \texttt{rateLimiter} · \texttt{errorHandler}};

  % ── Tier 2 : Business Logic / Controller Layer ────────────────────
  \node[layer] (controllers) at (0, 3.5) {Tier 2 — Business Logic / Controller Layer};
  \node[module] at (-3.5, 2.6) {\texttt{authController}};
  \node[module] at (-1.1, 2.6) {\texttt{assetController}};
  \node[module] at (1.3,  2.6) {\texttt{orderController}};
  \node[module] at (3.7,  2.6) {\texttt{userController}};

  % ── Tier 3 : Data Access / Repository Layer ───────────────────────
  \node[layer] (repos) at (0, 1.8) {Tier 3 — Data Access / Repository Layer};
  \node[module] at (-3.5, 0.9) {\texttt{userRepository}};
  \node[module] at (-1.1, 0.9) {\texttt{assetRepository}};
  \node[module] at (1.3,  0.9) {\texttt{orderRepository}};
  \node[module] at (3.7,  0.9) {\texttt{cartRepository}};

  % ── Database ──────────────────────────────────────────────────────
  \node[cylinder, draw=darkGray, fill=gray!12, shape border rotate=90,
        minimum width=4cm, minimum height=0.9cm, align=center,
        font=\footnotesize\bfseries] at (0, 0) {PostgreSQL Database};

  % ── Cross-cutting Services ────────────────────────────────────────
  \node[font=\footnotesize\bfseries\color{orange!80!black}]
        at (6.7, 3.5) {Cross-cutting};
  \node[cross] at (6.7, 2.9) {\texttt{storageService}};
  \node[cross] at (6.7, 2.0) {\texttt{auditService}};
  \node[cross] at (6.7, 1.1) {\texttt{aiService}};
  \node[cross] at (6.7, 0.2) {\texttt{tokenBlacklist}};

  % ── Arrows between tiers ─────────────────────────────────────────
  \draw[arr] (routes) -- node[lbl] {delegates} (controllers);
  \draw[arr] (controllers) -- node[lbl] {queries} (repos);
  \draw[arr] (repos.south) -- node[lbl] {SQL} (0, 0.45);

\end{tikzpicture}
}
\caption{Backend Layered Architecture --- Routes → Controllers → Repositories}
\label{fig:arch_backend}
\end{figure}

\section*{Conclusion}
This chapter identified all system actors, detailed functional and non-functional requirements, and established the architectural foundation. The platform follows a three-tier distributed architecture with the Repository pattern on the backend. The following chapters each cover one sprint or pair of sprints of the development cycle.
